\documentclass[../../../uem_mat_mei_q.tex]{subfiles}

\begin{document}
    2つの関数$f(x)=x^3+3tx^2+12(t-1)x$，$g(x)=-f(x)$について，%
    座標平面上の2曲線$y=f(x)$，$y=g(x)$を考える．ただし，$t$は定数である．%
    また，関数$f(x)$が極大値と極小値をもち，極値をとるときの$x$の値を%
    $\alpha$，$\beta$（ただし，$\alpha<\beta$）とする．

    以下の問に答えなさい．設問\,\ref{q:uem_mat_mei_2025_bus_03_q:1}，%
    \ref{q:uem_mat_mei_2025_bus_03_q:2}\,は空欄内の各文字に当てはまる数字を%
    所定の解答欄にマークしなさい．設問\,\ref{q:uem_mat_mei_2025_bus_03_q:3}，%
    \ref{q:uem_mat_mei_2025_bus_03_q:4}，\ref{q:uem_mat_mei_2025_bus_03_q:5}\,は%
    裏面の所定の欄に解答のみ書きなさい．ただし，分数はすべて既約分数にしなさい．

    \begin{enumerate}
        \item%
            $\alpha+\beta=-\:\myboxb{ユ}\:t$，$\alpha\beta=\:\myboxb{ヨ}\:t-\myboxb{ラ}$\,，\\%
            $\beta-\alpha=\left|\,\myboxb{リ}\:t-\:\myboxb{ル}\,\right|$ である．
            \label{q:uem_mat_mei_2025_bus_03_q:1}
        \item%
            2曲線$y=f(x)$，$y=g(x)$の共有点が3個となるとき，原点以外の共有点の$x$座標を%
            $\gamma$，$\delta$とすると，\\
            $\gamma+\delta=\:\myboxb{レ}\:t$，%
            $\gamma\delta=\:\myboxb{ロワ}\:t-\:\myboxb{ヲン}$ である．
            \label{q:uem_mat_mei_2025_bus_03_q:2}
        \item%
            2曲線$y=f(x)$，$y=g(x)$の共有点が2個となる$t$の値をすべて書きなさい．
            \label{q:uem_mat_mei_2025_bus_03_q:3}
        \item %
            2曲線$y=f(x)$，$y=g(x)$の共有点が3個となるとき，%
            4点$\mathrm{A}(\alpha,\:f(\alpha))$，$\mathrm{B}(\alpha,\:g(\alpha))$，%
            $\mathrm{C}(\beta,\:f(\beta))$，$\mathrm{D}(\beta,\:g(\beta))$を頂点とする%
            四角形ABCDの面積を$S$とする．$S$を$t$を用いて表した式を書きなさい．
            \label{q:uem_mat_mei_2025_bus_03_q:4}
        \item%
            \ref{q:uem_mat_mei_2025_bus_03_q:4}\,の条件を満たし，かつ，%
            $\beta-\alpha\geqq3$となるとき，四角形ABCDの面積が最小となる$t$の値と%
            そのときの面積$S$を求めなさい．
            \label{q:uem_mat_mei_2025_bus_03_q:5}
    \end{enumerate}
\end{document}
